\documentclass[../main.tex]{subfiles}

\begin{document}

\section{Sensitivity Analysis}\label{sec:sensitivity}

To see how fuel consumption depends on the imposed flight time, I re-solve the NLP at $t_f \in \{0.95,\, 1.00,\, 1.05\}\cdot 38$~s, i.e., at $36.10$, $38.00$, and $39.90$~s. Each solve uses the linear-interpolation initial guess of Section~\ref{sec:initial_guess} (no warm-start across runs). The shortest and nominal runs ($t_f = 36.10$~s and $38.00$~s) stop at the \texttt{MaxIterations} cap with first-order optimality of order $10^{-3}$--$10^{-4}$ in the non-dim cost; the longest run ($t_f = 39.90$~s) meets the step-size convergence criterion. The two capped runs still satisfy all constraints to within the \texttt{ConstraintTolerance} of $10^{-6}$, so I report them as engineering-quality solutions rather than fully optimised ones.

Table~\ref{tab:hm2_sweep} reports the converged final masses and propellant consumption.

\begin{table}[H]
\centering
\caption{Sensitivity sweep on flight time.}
\label{tab:hm2_sweep}
\begin{tabular}{ccc}
\toprule
$t_f$ [s] & $m_f$ [kg] & Fuel consumed [kg] \\
\midrule
36.10 & 1406.33 & 593.67 \\
38.00 & 1403.20 & 596.80 \\
39.90 & 1398.82 & 601.18 \\
\bottomrule
\end{tabular}
\end{table}

Fuel use grows with $t_f$ across the swept window. The reason is simple: a longer descent keeps the vehicle in the gravity field longer, so more thrust goes into counteracting gravity (gravity loss). Within the $\pm 5\%$ window around the nominal $38$~s the variation is small ($\approx 7.5$~kg, or about $1.3\%$ of the propellant), so the nominal $t_f$ sits near a flat region of the cost---which fits the fact that the mission is dominated by horizontal divert and not by hover time.

The change in flight time is absorbed almost entirely by the coast arc. Table~\ref{tab:hm2_switch} shows this through the switching times, computed as the crossings of the thrust magnitude through $0.5\,T_{\max}$: the first burn barely moves across the sweep, while the start of the terminal burn shifts to track the imposed $t_f$.

\begin{table}[H]
\centering
\caption{Thrust switching structure across the sensitivity sweep.}
\label{tab:hm2_switch}
\begin{tabular}{cccc}
\toprule
$t_f$ [s] & Burn 1 ends [s] & Burn 2 starts [s] & Coast length [s] \\
\midrule
36.10 & 13.91 & 31.32 & 17.41 \\
38.00 & 14.04 & 33.12 & 19.09 \\
39.90 & 14.08 & 35.04 & 20.95 \\
\bottomrule
\end{tabular}
\end{table}

The monotonic trend suggests that the fuel-optimal solution sits at the smallest $t_f$ the thrust-magnitude bound allows. A free-time formulation, where $t_f$ becomes a decision variable, would find this minimum on its own; this is one of the extensions listed in Section~\ref{sec:conclusions}.

The trajectory, thrust, mass, and glide-slope plots in Section~\ref{sec:results} overlay the three sensitivity runs on common axes for direct comparison. The solutions look much alike: all three show the max--coast--max thrust structure, all three stay inside the glide-slope corridor, and the trajectory shapes differ only slightly, with the extra flight time spent coasting at altitude before the terminal burn.

\end{document}
